\subsection{Conclusions}


In this paper we have studied the differential source-count distribution of gamma-ray sources below the \fermi  threshold for source identification. We  adopted a convolutional neural network methodology devised to deal with spherical sky-maps and able to extract information related to small scale fluctuations. The technique is based on the EfficientNet V2M model, applied to a patched version of the 2-dimensional spherical sky-map projection of the gamma-ray emission. The neural network has been trained on 900k synthetic maps and validated on 100k additional maps. These synthetic maps have been generated from the flux of random selections of source-counts distributions, to which galactic foreground models and an isotropic component have been added. After training and validation, the neural network has been applied to the \fermi sky-map.

We concentrated our analysis on the photon energy interval $(1,10)$ GeV and to 14 years of data. The main result is shown in Fig. \ref{fig:prediction-fermi-12y}: the reconstructed source count distribution exhibits a $\dnds \sim S^{-2}$ behaviour over almost four orders of magnitude in flux in the range  
$[5 \cdot 10^{-12} , 1 \cdot 10^{-8}]$ cm$^{-2}$ s$^{-1}$, which includes the unresolved regime range 
$[5 \cdot 10^{-12}, 2 \cdot 10^{10}]$.
In the regime where \fermi has sensitivity to individually resolve gamma-ray sources, Fig. \ref{fig:prediction-fermi-12y} shows that the neural network fairly reproduces the observations, giving confidence on its reliability and strength.
While in the unresolved regime the result is in agreement with previous studies performed with different methodologies.
We have also further validated the result performing several test of stability which have confirmed that  our findings is stable and robust.

 The methodology presented here provides a proof of principle for the adoption of a CNN to the reconstruction of the source-count distribution of the extra-galactic sky. 
With respect to other methods to extract the source-count distribution in the unresolved regime, like the 1-point PDF one, the use of a CNN avoids the need to calculate complicated and numerically demanding likelihoods.
 Possible future applications includes the extension to multiple energy ranges and energy correlations, and the investigation of features in the $dN/dS$ which might indicate the presence of exotic components like dark matter.
A further aspect which would be interesting to investigate is to compare different approaches to machine learning on spherical domains. In particular, it might be interesting to employ a graph convolutional neural network as suggested in \cite{DBLP:journals/corr/abs-2012-15000}, in order to preserve spatial relations between distant pixels, and rotational invariance.
Finally, recent developments in the field of simulation based inference, such as our case, suggest that a promising approach to estimate the complex posterior distribution of the objective parameters are the so-called neural likelihood ratio estimation techniques \cite{Durkan2020}, and in particular Truncated Marginal Neural Ratio Estimation \cite{Miller2021arxiv,Miller2021,AnauMontel:2022ppb}. Since these techniques are tailored for simulation based problems, we hypothesise that it would be possible to achieve similar performance to our current neural network architecture with an even simpler structure and possibly less training samples. 
%This would also make the extension to a multiple energy bins analysis a much more feasible and less computationally expensive task.  
 
 

\label{sec:dnds:conclusions}

